\documentclass[10pt]{article}
% Math Packages
\usepackage{amsmath, mathtools}
\usepackage{amssymb}
\usepackage{amsthm}
\usepackage{amsfonts}
\usepackage{bbm}
\usepackage{breqn}
\usepackage[margin=1in]{geometry}
\usepackage{graphicx}
\usepackage{tikz}
\usetikzlibrary{arrows.meta}
\usetikzlibrary{calc}
\usepackage{forest}
\usepackage{tikz-qtree}
\graphicspath{ {./images/} }
\usepackage{hyperref}
\usepackage[capitalize]{cleveref}
\usepackage[shortlabels]{enumitem}
\usetikzlibrary{arrows,matrix,positioning}
\usepackage{multicol}

\usepackage{listings}
\usepackage{xcolor}

\crefname{problem}{Problem}{Problems}
\Crefname{problem}{Problem}{Problems}

\lstset{
  language=R,
  basicstyle=\ttfamily\footnotesize,
  keywordstyle=\color{blue},
  commentstyle=\color{gray},
  stringstyle=\color{red!70!black},
  framesep=2pt,
  framexleftmargin=0pt, 
  frame=single,
  breaklines=false,
  showstringspaces=false
}
% for the pipe symbol
\usepackage[T1]{fontenc}

% Colorful Notes
\usepackage{color} \definecolor{Red}{rgb}{1,0,0} \definecolor{Blue}{rgb}{0,0,1}
\definecolor{Purple}{rgb}{.5,0,.5} \def\red{\color{Red}} \def\blue{\color{Blue}}
\def\gray{\color{gray}} \def\purple{\color{Purple}}
\newcommand{\rnote}[1]{{\red [#1]}} % \rnote{foo} gives '[foo]' in red
\newcommand{\pnote}[1]{{\purple [#1]}} % \pnote{foo} gives '[foo]' in purple
\newcommand{\bnote}[1]{{\blue #1}} % \bnote{foo} gives 'foo' in blue
\newcommand{\gnote}[1]{{\gray #1}} % \gnote{foo} gives 'foo' in gray
\newcommand{\Max}[1]{{\purple [#1]}} % \bnote{foo} then 'foo' is blue

% Claim numbering (the counter restarts after each proof environment)
\newcounter{claimcount}
\setcounter{claimcount}{0}
\newenvironment{claim}{\refstepcounter{claimcount}\par\addvspace{\medskipamount}\noindent\textbf{Claim \arabic{claimcount}:}}{}
\usepackage{etoolbox}
\AtBeginEnvironment{proof}{\setcounter{claimcount}{0}}
\newenvironment{claimproof}{\par\addvspace{\medskipamount}\noindent\textit{Proof of Claim  \arabic{claimcount}.}}{\hfill\ensuremath{\qedsymbol} \tiny{Claim}

  \medskip}
% Add claim support to cleverref
\crefname{claimcount}{Claim}{Claims}

% Math Environments
\newtheorem{theorem}{Theorem}
\newtheorem{assumption}[theorem]{Assumption}
\newtheorem{lemma}[theorem]{Lemma}
\newtheorem{proposition}[theorem]{Proposition}
\newtheorem{corollary}[theorem]{Corollary}
\newtheorem{question}[theorem]{Question}
\theoremstyle{definition}
\newtheorem*{definition}{Definition}
\newtheorem{remark}[theorem]{Remark}
\newtheorem{example}[theorem]{Example}
\newtheorem{notation}[theorem]{Notation}
\newtheorem{problem}[theorem]{Problem}

% Matrices and Column Vectors. 
\usepackage{stackengine}
\setstackgap{L}{1.0\normalbaselineskip}
\usepackage{tabstackengine}
\setstackEOL{;}% row separator
\setstackTAB{,}% column separator
\setstacktabbedgap{1ex}% inter-column gap
\setstackgap{L}{1.5\normalbaselineskip}% inter-row baselineskip
\let\nmatrix\bracketMatrixstack  %Usage: \nmatrix{1,2,3\4,5,6}
\newcommand\cv[1]{\setstackEOL{,}\bracketMatrixstack{#1}} %usage: \cv{1,2,3}

% Custom Math Coqmmands
\newcommand{\vt}{\vskip 5mm} % vertical space
\newcommand{\fl}{\noindent\textbf} % first line
\newcommand{\Fl}{\vt\noindent\textbf} % first line with space above
\newcommand{\norm}[1]{\left\lVert#1\right\rVert} % norm
\newcommand{\pnorm}[1]{\left\lVert#1\right\rVert_p} % p-norm
\newcommand{\qnorm}[1]{\left\lVert#1\right\rVert_q} % q-norm
\newcommand{\1}[1]{\textbf{1}_{\left[#1\right]}} % indicator function 
\def\limn{\lim_{n\to\infty}} % shortcut for lim as n-> infinity
\def\sumn{\sum_{n=1}^{\infty}} % shortcut for sum from n=1 to infinity
\def\sumkn{\sum_{k=1}^{n}} % shortcut for sum from k=1 to n
\def\sumin{\sum_{i=1}^{n}} % shortcut for sum from i=1 to n
\def\SAs{\sigma\text{-algebras}} % shortcut for $\sigma$-algebras
\def\SA{\sigma\text{-algebra}} % shortcut for $\sigma$-algebra
\def\Ft{\mathcal{F}_t} % time-indexed sigma-algebra (t)
\def\Fs{\mathcal{F}_s} % time-indexed sigma-algebra (s)
\def\F{\mathcal{F}} % sigma-algebra
\def\G{\mathcal{G}} % sigma-algebra
\def\R{\mathbb{R}} % Real numbers
\def\N{\mathbb{N}} % Natural numbers
\def\Z{\mathbb{Z}} % Integers
\def\E{\mathbb{E}} % Expectation
\def\P{\mathbb{P}} % Probability
\def\Q{\mathbb{Q}} % Q probability
\def\dist{\text{dist}} %Text 'dist' for things like 'dist(x,y)'
\newcommand{\indep}{\perp \!\!\! \perp}  %independence symbol
\def\Var{\mathrm{Var}} % Variance
\def\tr{\mathrm{tr}} % trace

% Brackets and Parentheses
\def\[{\left [}
    \def\]{\right ]}
% \def\({\left (}
%   \def\){\right )}

\usepackage{color}
\definecolor{Red}{rgb}{1,0,0}
\definecolor{Blue}{rgb}{0,0,1}
\definecolor{Purple}{rgb}{.5,0,.5}
\def\red{\color{Red}}
\def\blue{\color{Blue}}
\def\gray{\color{gray}}
\def\purple{\color{Purple}}
\definecolor{RoyalBlue}{cmyk}{1, 0.50, 0, 0}
\newcommand{\dempfcolor}[1]{{\color{RoyalBlue}#1}} 
\newcommand{\demph}[1]{\textcolor{RoyalBlue}{\textbf{\slshape #1}}} % Slanted

\usepackage{emoji}
% comment exactly one of the following line to show / hide the solutions
% \newcommand{\solution}[1]{{\purple #1}} % uncomment to show the solutions
\newcommand{\solution}[1]{} % uncomment to hide the solutions




\begin{document}
\begin{center}
  \section*{Math 472: Homework 08}
  \textit{Due Tuesday, March 31}
\end{center}


\begin{problem}[Exercise 9.82]
  Let $Y_{1},\ldots,Y_{n}$ be a random sample from the density function given
  by
  \begin{equation*}
    f(y\mid \theta)  = \left\{
      \begin{array}{l@{\quad:\quad}l}
        \left( \frac{1}{\theta} \right)r y^{r-1}e^{-y^{r-1}}e^{-y^{r}/\theta}
        & \theta>0, y>0\\
        0& \text{elsewhere.}\end{array}\right.
  \end{equation*}
  where $r>0$ is a known positive constant.
  \begin{enumerate}[(a)]
    \item Find a sufficient statistic for $\theta$.
    \item Find the MLE of $\theta$. 
  \end{enumerate}
\end{problem}


\begin{problem}
  Let $X_{1},\ldots,X_{n}$ be a random sample from the pdf
  \begin{equation*}
    f(x\mid \theta) = \theta x^{-2}, \quad 0<\theta \leq x <\infty.
  \end{equation*}
  \begin{enumerate}[(a)]
    \item Find a suffificent statistic for $\theta$. 
    \item Find the MLE of $\theta$. 
  \end{enumerate}
\end{problem}


\begin{problem}[Shifted exponential]
  A water supply becomes contaiminated at some unknown time $\mu$.
  
  Suppose there are a total of $n$ sensors in the water. After the
  contamination, these sensors don't immediately start testing positive, but
  more realistically do so with with some lag times (owing to sensor location,
  sensitivity, and uneven contaminant concentrations).

  For each $i\in [n]$,
  let $X_{i}$ be the first time that sensor $i$ tests positive for
  contamination. Assume that $X_{1},\ldots,X_{n}$ are independent and each
  have distribution given by the pdf
  \begin{equation*}
    f(x\mid \theta) = e^{-(x-\mu)}\mathbf{1}_{\left[\mu \leq x<\infty\right]}.
  \end{equation*}
  \begin{enumerate}[(a)]
    \item Show that $X_{(1)}:=\min(X_{1},\ldots,X_{n})$ is a sufficient
    statistic for $\mu$. 
    \item Find the MLE of $\mu$.
  \end{enumerate}
\end{problem}


\begin{problem}[Exercise 10.2]
  \label{problem:exercise-10.2}
  An experimenter has prepared a drug dosage level that she claims will induce
  sleep for 80\% of people suffering from insomnia. After examining the
  dosage, we feel her claims are inflated. In an attempt to disprove her
  claim, we administer the drug to 20 insomniacs and observe $Y$, the number
  for whom the drug dose induces sleep. We wish to test the hypothesis
  $H_{0}:p=.8$ vs the alternative $H_{a}:p<.8$. Assume that the rejection
  region $\left\{y \leq 12\right\}$ is used.
  \begin{enumerate}[(a)]
    \item In terms of this problem, what is the type I error?
    \item Find $\alpha$.
    \item In terms of this problem, what is the type II error?
    \item Find $\beta$ when $p=.6$.
    \item Find $\beta$ when $p=.4$.
  \end{enumerate}
\end{problem}


\begin{problem}[Exercise 10.3]
  Refer to \cref{problem:exercise-10.2}.
  \begin{enumerate}[(a)]
    \item Find the rejection region of the for $\left\{y \leq c\right\}$ so
    that $\alpha\approx .01$.
    \item For the the rejection region in part (a), find $\beta$  when $p=.6$.
    \item For the rejection region in part (a), find $\beta$ when $p=.4$.
  \end{enumerate}
\end{problem}


\begin{problem}[Exercise 10.19]
  \label{problem:exercise-10.19}
  The output voltage for an electric circuit is specified to be 130. A sample
  of 40 independent readings on the voltage for this circuit gave a sample
  mean 128.6 and standard deviation 2.1. Using the large-sample z-test method
  described in section 10.3 of the textbook, test the hypothesis that the average
  output voltage is 130 against the alternative that it is less than 130. Use
  a test with level .05.
\end{problem}


\begin{problem}[Exercise 10.37]
  Refer to \cref{problem:exercise-10.19}. If the voltage falls as low as 128,
  serious consequences may result. For testing $H_{0}: \mu=130$ versus
  $H_{a}:\mu=128$, find the probability of a type II error, $\beta$ for the
  rejection region used in \cref{problem:exercise-10.19}
\end{problem}


\begin{problem}[Exercise 10.20]
  \label{problem:exercise-10.20}
  The Rockwell hardness index for steel is determined by pressing a diamond
  point into the steel and measuring the depth of penetration. For 50
  specimens of an alloy of steel, the Rockwell hardness index averaged 62 with
  standard deviation 8. The manufacturer claims that this alloy has an average
  hardness index of at least 64. Is there sufficient evidence to refute the
  manufacturer's claim at the 1\% significance level? Use the large-sample
  z-test method described in section 10.3 of the textbook.
\end{problem}


\begin{problem}[Exercise 10.38]
  Refer to problem \cref{problem:exercise-10.20}. The steel is sufficiently
  hard to meet usage requirements if the mean Rockwell hardness measure does
  not drop below 60. Using the rejection region found in 
  \cref{problem:exercise-10.20}, find $\beta$ for the specific alternative
  $\mu_{a} = 60$.
\end{problem}


\begin{problem}[Salk Polio Field Trial (1954)]
  \label{ex:salk-polio-field-trial-1954}
  The results of a 1954 polio vaccine trial in $n=401,974$ elementary school
  students gave the following results:
  \begin{equation}\label{eq:table-of-observed-values}
    \begin{tabular}{l|cc}
      & \textbf{Polio} & \textbf{No Polio} \\
      \hline
      \textbf{Vaccinated} & 200,712          & 33            \\
      \textbf{Unvaccinated}  & 201,114           & 115           
    \end{tabular} 
  \end{equation}
  We wish to know if there is a relationship between vaccination status and
  the incidience of paralytic polio. The null hypothesis is:
  \begin{equation*}
    H_{0}= \text{vaccination and polio are independent}
  \end{equation*}
  and the alternative hypothesis is
  \begin{equation*}
    H_{a} = \text{ there is a relationship between vaccination and polio.}
  \end{equation*}
  Conduct a hypothesis test using the chi-squared statistic to test these hypotheses.
\end{problem}

\end{document}